\section{Discussion}
\label{sec:discussion}

\subsection{Planning is not feedback}
\label{sec:planning-not-feedback}

Each rung of the ladder measures the value of a different capability:

\begin{center}
\begin{tabular}{ll}
\toprule
Rung & Capability measured \\
\midrule
constant $\rightarrow$ one-step   & value of immediate per-state selection \\
one-step $\rightarrow$ committed  & value of multi-step lookahead \\
committed $\rightarrow$ replanning & value of revising the plan during execution \\
\bottomrule
\end{tabular}
\end{center}

\textbf{We do not decompose these values arithmetically.} The median of paired
differences is not linear. In our measurements \ctrl{committed} is $+2.090$ and
replanning is $+1.690$ against the tuned constant, yet the directly measured
$\text{replanning} - \ctrl{committed}$ is $+0.010$. By spec the values are
$+0.472 / -0.019 / +0.009 / +0.000$, so the two planners are effectively tied and
the marginal gap arises because the pooled medians fall on different instances.
\emph{Comparisons are made only with direct paired statistics}
(Section~\ref{sec:no-subtraction}, Figure~\ref{fig:planning-vs-feedback}).

\textbf{Q1} (does a good sequence exist?) is supported; \textbf{Q2} (is it worth
revising during execution?) is not supported under these conditions. On a
deterministic objective the planner's internal model is exact, so the plan formed at
the initial state is already the optimal prediction and replanning gains no new
information.

\subsection{Model error raises the risk of long-horizon plans}
\label{sec:model-error}

The exploratory micro-neural results summarize as follows.

\begin{quote}
Long-horizon plans are valuable when the internal predictive model is accurate; when
the model is wrong, the risk of a stale plan grows.
\end{quote}

In the minibatch regimes the rejection rate of the committed plan rose from $0.00$ to
$0.66$--$0.79$ and terminal improvement fell below the tuned constant. Replanning
reduced that loss.

\textbf{Yet under the same conditions replanning is worse than one-step control.}
$\text{replanning} - \text{one-step}$ is $-1.111$ and $-0.716$ in the two minibatch
conditions, and $\text{replanning} - \text{tuned constant}$ is $-0.900$ and $-0.277$.
So one must not read the positive $\ctrl{committed} \rightarrow \text{replanning}$
value as evidence for a learned feedback policy. \emph{That positive value is
stale-plan avoidance, not an advantage over inexpensive myopic control.}

At $n = 3$ per regime we make no claim about magnitude.

\subsection{Why the benchmark audit was necessary}
\label{sec:audit-why}

\begin{quote}
Our benchmark audit was necessary because nominal numerical ceilings and task labels
did not reliably identify whether a problem could distinguish controllers.
\end{quote}

These are the problems that actually reversed conclusions in this project:

\begin{itemize}
  \item joint saturation at the numerical floor (two of three specs in the initial
        dev subset);
  \item one-sided saturation, where dropping pairs removed the good results;
  \item an attainable local minimum of Rosenbrock, which passed the eligibility
        check;
  \item identical seeds producing identical instances, so a nominal $n = 3$ was
        $n = 1$;
  \item the difference between a ceiling based on the global optimum and one based on
        the starting-point basin;
  \item the possibility that a single reference solver fails (L-BFGS did not
        converge at $\kappa = 10^6$).
\end{itemize}

Both of our initial rules --- ``judge saturation from the task name'' and ``compute
the ceiling from the numerical floor'' --- were wrong. That is why we introduced the
reference-solver panel and the attainable ceiling.

\textbf{We are not proposing a general benchmark framework.} This is an empirical
procedure that corrected this project's benchmark, and we intend it to be read as a
checklist to consult when designing similar controller comparisons.

\subsection{What the cheap method misses}
\label{sec:cheap}

\ctrl{onestep} spends $7.9$ times the budget on search --- $1/164$ of the planner
--- and gains $+1.155$~nat over the tuned constant. The planner gains $+1.690$~nat
but spends $1{,}294$ times the budget. The directly measured increment of
$+0.456$~nat is statistically robust, but this cost gap is its price. From a
practical standpoint this is the most direct implication of the present results.
